\documentclass[10pt]{extarticle}
\usepackage[a4paper,margin=1.5cm]{geometry}
\usepackage[T1]{fontenc}
\usepackage{amsmath,amssymb}
\usepackage{graphicx}
\usepackage[export]{adjustbox}
\usepackage{multicol}
\usepackage{enumitem}
\usepackage{booktabs,array,tabularx}
\setlength{\columnsep}{0.7cm}
\setlength{\parindent}{0pt}
\setlength{\parskip}{3pt}
\pagestyle{empty}
\makeatletter
% figure registry (figdefs.tex is generated by the extract step)
\newcommand{\qfigdef}[4]{\@namedef{qf@file@#1}{#2}\@namedef{qf@w@#1}{#3}\@namedef{qf@wide@#1}{#4}}
\newcommand{\qf@check}[1]{\@ifundefined{qf@file@#1}{\PackageError{score2}{Unknown figure id #1}{}}{}}
\newcommand{\qf@img}[4][]{\includegraphics[width=#2\dimexpr\@nameuse{qf@w@#4}pt\relax,max width=#3,max height=0.45\textheight,#1]{\@nameuse{qf@file@#4}}}
% Full-width questions: finalize wraps every question that contains a wide figure in \qfullbegin...\qfullend,
% so the WHOLE question (stem, figures, options) is set across the page in reading order, then the two columns
% resume. The question's head - its number, stem and first figure (\qfullsplit marks the end of the head) -
% is kept on one page, so a number is never left at a page foot with its figure overleaf; the rest may break
% across pages normally, so no big gaps are left. Without a \qfullsplit, the whole question is kept together
% if it is under 45% of a page. Never write these yourself.
\newif\ifqf@full
\newif\ifqf@boxed
\newsavebox{\qf@box}
\newcommand{\qf@open}{\global\qf@fulltrue\par\global\qf@boxedtrue
  \setbox\qf@box\vbox\bgroup\hsize\textwidth\linewidth\textwidth}
\newcommand{\qf@emit}[1]{\par\egroup\global\qf@boxedfalse
  \ifdim\dimexpr\ht\qf@box+\dp\qf@box\relax<#1\textheight\noindent\usebox{\qf@box}\par\else\unvbox\qf@box\fi}
\newcommand{\qfullsplit}{\ifqf@boxed\qf@emit{0.85}\fi}
\newcommand{\qf@close}{\ifqf@boxed\qf@emit{0.45}\fi\global\qf@fullfalse}
\newcommand{\qfullbegin}{\end{multicols}\qf@open}
\newcommand{\qfullend}{\qf@close\begin{multicols}{2}}
\newcommand{\qfullstart}{\qf@open}
\newcommand{\qfulldone}{\qf@close}
% \qfig[scale]{id}: block figure, centred, directly where it is written (size set by the tool).
% Fallback if a wide figure is not inside a full-width question: break out of the columns in place - no floats.
\newcommand{\qfig}[2][1]{\qf@check{#2}%
  \ifqf@full\qf@blk{#1}{#2}\else
  \ifnum\@nameuse{qf@wide@#2}=1 \expandafter\@firstoftwo\else\expandafter\@secondoftwo\fi
  {\end{multicols}\begin{center}\qf@img{#1}{\textwidth}{#2}\end{center}\begin{multicols}{2}}%
  {\qf@blk{#1}{#2}}\fi}
\newcommand{\qf@blk}[2]{\par\smallskip{\centering\qf@img{#1}{\linewidth}{#2}\par}\smallskip}
% \qfigi[scale]{id}: inline figure (inside option lists / tables / minipages); never breaks columns.
\newcommand{\qfigi}[2][1]{\qf@check{#2}\qf@img[valign=c]{#1}{\linewidth}{#2}}
\makeatother
% structure macros - lint reads these to check completeness and figure placement
\newcommand{\papersec}[1]{\par\bigskip{\centering\Large\bfseries PAPER-#1\par}\medskip}
\newcommand{\partsec}[2]{\par\medskip\noindent{\large\bfseries PART-#1: #2}\par\smallskip}
\newcommand{\sect}[1]{\par\smallskip\noindent\textbf{#1}\par}
\newcommand{\qn}[1]{\par\smallskip\noindent\textbf{#1.}~}
\newcommand{\paraq}[2]{\par\smallskip\noindent\textbf{Paragraph for Questions #1 and #2}\par}
\newenvironment{qoptions}{\begin{itemize}[leftmargin=1.8em,itemsep=1pt,topsep=2pt,label={}]}{\end{itemize}}
\newcommand{\opt}[2]{\item[(#1)] #2}
% figure registry (generated by score2_tools.py finalize)
\qfigdef{p1_phy_s1_q02_1}{Score Advanced - Test 12 (11 May 2023)-img/p1_phy_s1_q02_1.png}{103.9}{0}
\qfigdef{p1_phy_s1_q03_1}{Score Advanced - Test 12 (11 May 2023)-img/p1_phy_s1_q03_1.png}{81.3}{0}
\qfigdef{p1_phy_s1_q04_1}{Score Advanced - Test 12 (11 May 2023)-img/p1_phy_s1_q04_1.png}{79.3}{0}
\qfigdef{p1_phy_s2_q07_1}{Score Advanced - Test 12 (11 May 2023)-img/p1_phy_s2_q07_1.png}{82.0}{0}
\qfigdef{p1_phy_s2_q08_1}{Score Advanced - Test 12 (11 May 2023)-img/p1_phy_s2_q08_1.png}{83.4}{0}
\qfigdef{p1_phy_s2_q09_1}{Score Advanced - Test 12 (11 May 2023)-img/p1_phy_s2_q09_1.png}{133.8}{0}
\qfigdef{p1_phy_s3_q02_1}{Score Advanced - Test 12 (11 May 2023)-img/p1_phy_s3_q02_1.png}{123.7}{0}
\qfigdef{p1_phy_s3_q03_1}{Score Advanced - Test 12 (11 May 2023)-img/p1_phy_s3_q03_1.png}{107.8}{0}
\qfigdef{p1_phy_s3_q04_1}{Score Advanced - Test 12 (11 May 2023)-img/p1_phy_s3_q04_1.png}{91.6}{0}
\qfigdef{p1_phy_s3_q05_1}{Score Advanced - Test 12 (11 May 2023)-img/p1_phy_s3_q05_1.png}{115.8}{0}
\qfigdef{p1_phy_s3_q06_1}{Score Advanced - Test 12 (11 May 2023)-img/p1_phy_s3_q06_1.png}{135.1}{0}
\qfigdef{p1_phy_s3_q08_1}{Score Advanced - Test 12 (11 May 2023)-img/p1_phy_s3_q08_1.png}{105.6}{0}
\qfigdef{p1_chem_s1_q01_1}{Score Advanced - Test 12 (11 May 2023)-img/p1_chem_s1_q01_1.png}{212.6}{0}
\qfigdef{p1_chem_s1_q02_1}{Score Advanced - Test 12 (11 May 2023)-img/p1_chem_s1_q02_1.png}{271.3}{1}
\qfigdef{p1_chem_s2_q06_1}{Score Advanced - Test 12 (11 May 2023)-img/p1_chem_s2_q06_1.png}{254.0}{1}
\qfigdef{p1_chem_s3_q01_1}{Score Advanced - Test 12 (11 May 2023)-img/p1_chem_s3_q01_1.png}{47.7}{0}
\qfigdef{p1_chem_s3_q05_1}{Score Advanced - Test 12 (11 May 2023)-img/p1_chem_s3_q05_1.png}{243.6}{0}
\qfigdef{p1_chem_s3_q06_1}{Score Advanced - Test 12 (11 May 2023)-img/p1_chem_s3_q06_1.png}{147.9}{0}
\raggedcolumns
\makeatletter\@beginparpenalty=10000\makeatother
\newcolumntype{L}[1]{>{\raggedright\arraybackslash}p{#1}}
\begin{document}
\begin{center}{\Large\bfseries Score Advanced --- Test 12 (11 May 2023)}\end{center}
\begin{multicols}{2}
\raggedcolumns
\papersec{1}
\partsec{1}{PHYSICS}
\sect{SECTION-I (i)}
\qn{1} At the initial moment three point A, B and C are on a horizontal plane along a straight line such that AB = BC. Point A begins to move upward with a constant velocity `v' and point C downward without any initial velocity at a constant acceleration `a'. If the points being to move simultaneously, then the initial velocity and acceleration of point B for all the three particles to be constantly on same straight line must be
\begin{qoptions}
\opt{A}{$\dfrac{v}{2}$ upwards, $\dfrac{a}{2}$ downwards}
\opt{B}{$\dfrac{v}{2}$ upwards, $\dfrac{a}{2}$ upwards}
\opt{C}{$\dfrac{v}{2}$ downwards, $\dfrac{a}{2}$ downwards}
\opt{D}{$\dfrac{v}{2}$ downwards, $\dfrac{a}{2}$ upwards}
\end{qoptions}
\qn{2} A wire (with some resistance) of mass m \& length $\ell$ can slide freely on a pair of parallel smooth horizontal rails. A uniform magnetic field B exist in the region in the direction perpendicular to the plane of the rails. The rails are connected at left end by a capacitor of capacity C. At t = 0, wire is given a velocity v$_0$ towards right side. Then\par
\qfig{p1_phy_s1_q02_1}
\begin{qoptions}
\opt{A}{Wire stops finally}
\opt{B}{Wire will oscillate}
\opt{C}{Wire will move with constant velocity after sufficient time.}
\opt{D}{Final K.E. of the rod must be equal to the energy stored in the capacitor}
\end{qoptions}
\qn{3} As shown, a piston chamber of cross section area A is filled with ideal gas. A sealed piston of mass m is right at the middle height of the cylinder at equilibrium. The friction force between the chamber wall and the piston can be ignored. The mass of the gas + chamber is M. The atmosphere pressure is P$_0$ . Now slowly pull the piston upwards, find the maximum value of M such that the chamber can be lifted off the ground. The temperature remains unchanged.\par
\qfig{p1_phy_s1_q03_1}
\begin{qoptions}
\opt{A}{M = $\dfrac{\text{P}_0\text{A} - 2\text{mg}}{2\text{g}}$}
\opt{B}{M = $\dfrac{\text{P}_0\text{A} - \text{mg}}{2\text{g}}$}
\opt{C}{M = $\dfrac{\text{P}_0\text{A} - \text{mg}}{\text{g}}$}
\opt{D}{M = $\dfrac{\text{P}_0\text{A}}{2\text{g}}$}
\end{qoptions}
\qn{4} Two bodies A and B have the same surface area and mass. The bodies have absolute temperature T$_\text{A}$ and T$_\text{B}$ , emissivities e$_\text{A}$ and e$_\text{B}$ , and specific heat capacities S$_\text{A}$ and S$_\text{B}$ . The intensity of radiation (E) is plotted against the wavelength ($\lambda$) of radiation for both bodies. Which of the following is possible?\par
\qfig{p1_phy_s1_q04_1}
\begin{qoptions}
\opt{A}{T$_\text{A}$ = T$_\text{B}$, e$_\text{A}$ = e$_\text{B}$}
\opt{B}{T$_\text{A}$ = T$_\text{B}$, e$_\text{A}$ $\neq$ e$_\text{B}$}
\opt{C}{T$_\text{A}$ $\neq$ T$_\text{B}$, e$_\text{A}$ = e$_\text{B}$,}
\opt{D}{T$_\text{A}$ $\neq$ T$_\text{B}$, e$_\text{A}$ = e$_\text{B}$}
\end{qoptions}
\sect{SECTION-I (ii)}
\qn{5} The pitch of screw gauge is 1 mm, and its circular scale is divided into 100 divisions. When nothing is put between the studs the zero of main scale is not seen but when circular scale is rotated by 450$^\text{o}$ the zero of main scale is just visible and the zero of main scale coincides with the zero of circular scale. When a glass plate is placed between the studs, the circular scale lies between 18$^\text{th}$ and19$^\text{th}$ division of main scale and circular scale reads 34 divisions. Then
\begin{qoptions}
\opt{A}{There is a positive zero error, and its magnitude is 1.25 mm.}
\opt{B}{There is a negative zero error, and its magnitude is 1.25 mm}
\opt{C}{The thickness of the glass plate is 19.59 mm.}
\opt{D}{The thickness of the glass plate is 17.09 mm}
\end{qoptions}
\qn{6} A transverse harmonic wave is propagating along a taut string. Tension in the string is 50 N and it's linear mass density is 0.02 kg/m. The string is driven by a 80 Hz oscillator tied to one end, oscillating with amplitude of 1 mm. The other end of string is terminated so that all the wave energy is absorbed and there is no reflection. From the given data, choose the \textbf{correct} option(s) (take $\pi^2$ =10)
\begin{qoptions}
\opt{A}{Power of the oscillator is 0.128 W}
\opt{B}{Average energy of the wave on a 1 m long segment of the string is 2.56 mJ}
\opt{C}{If tension is the string is quadrupled and power of oscillator remain same, then new amplitude of oscillation of string is $\dfrac{1}{\sqrt{2}}$ mm}
\opt{D}{Power of oscillator is 0.252 W}
\end{qoptions}
\qn{7} An isolated metallic spherical thin shell has given a charge Q as shown. The parameter r is the distance to be measured from center of sphere. Which of the following statement(s) is/are \textbf{CORRECT}.\par
\qfig{p1_phy_s2_q07_1}
\begin{qoptions}
\opt{A}{The energy stored in space R $<$ r $<$ 2R is $\dfrac{\text{Q}^2}{16\pi\varepsilon_0\text{R}}$}
\opt{B}{The energy stored in space R $<$ r $<$ 2R is $\dfrac{\text{Q}^2}{32\pi\varepsilon_0\text{R}}$}
\opt{C}{Energy stored in space 0 $<$ r $<$ $\dfrac{\text{R}}{2}$ is $\dfrac{\text{Q}^2}{16\pi\varepsilon_0\text{R}}$}
\opt{D}{If switch S is closed the heat generated in wire is $\dfrac{\text{Q}^2}{8\pi\varepsilon_0\text{R}}$}
\end{qoptions}
\qn{8} Magnetic field exist in a cylindrical region of radius `R' and centre `O' as shown in figure. A conducting rod PQ of length `a' is kept in the field as shown in figure. OPQR forms a square of side `a'. Magnetic field in the region is changing at the rate $\dfrac{\text{dB}}{\text{dt}}$ = $\alpha$($\alpha$ $>$ 0) then\par
\qfig{p1_phy_s2_q08_1}
\begin{qoptions}
\opt{A}{Line integral $\displaystyle\int_\text{P}^\text{Q}\vec{\text{E}}_\text{ind}.\vec{\text{dl}} = \dfrac{-\text{a}^2\alpha}{2}$}
\opt{B}{Potential difference between P \& Q is $\dfrac{\text{a}^2\alpha}{2}$}
\opt{C}{Line integral $\displaystyle\int_\text{P}^\text{Q}\vec{\text{E}}_\text{ind}.\vec{\text{dl}} = \dfrac{\text{a}^2\alpha}{2}$}
\opt{D}{Potential difference between Q \& R is $\dfrac{\text{a}^2\alpha}{2}$}
\end{qoptions}
\qn{9} As shown below, a point light source (S) is inside a glass block of refractive index n$_\text{G}$ = $\dfrac{3}{2}$, at a distance 2 mm from the optical axis and distance x = 3 cm from the flat surface. A thick plano-convex lens made of transparent plastic material of refractive index n$_\text{P}$ = 2, is placed near the source as shown. The thickness of the lens is d, and the radius of convex surface is R = 15 cm. The width of air gap between the glass block and the lens is `t'. The width of air gap t and thickness of lens d is adjusted so that light (paraxial rays) from point source forms always a parallel beam after passing through the lens. Select \textbf{CORRET} option(s).\par
\qfig{p1_phy_s2_q09_1}
\begin{qoptions}
\opt{A}{For d = 16cm, width of air gap t = 5 cm}
\opt{B}{For thickness d = 10 cm, width of air gap t = 8 cm}
\opt{C}{The maximum adjustable value of d = 26 cm.}
\opt{D}{The maximum adjustable value of d = 16 cm.}
\end{qoptions}
\qn{10} A $\longrightarrow$ B + C\par
A radioactive nucleus A at rest disintegrates into two nuclei B and C. Mass of B is 12 m and that of C is 4m. The Q-value of the reaction is, Q = $\dfrac{\text{h}^2}{24\text{m}\lambda^2}$. The energy liberated in reaction is completely imparted to the products (B and C) as kinetic energy.
\begin{qoptions}
\opt{A}{The de Broglie wavelength of B is 2$\lambda$}
\opt{B}{The de Broglie wavelength of C is 2$\lambda$}
\opt{C}{The mass of the nucleus A is 16m + $\dfrac{\text{h}^2}{24\text{mc}^2\lambda^2}$}
\opt{D}{The de Broglie wavelength of C is $\dfrac{2\lambda}{3}$}
\end{qoptions}
\sect{SECTION-III}
\qn{1} A horizontal stretched string, fixed at two ends, is vibrating in its 39$^\text{th}$ overtone according to the equation, y(x,t)=(0.01m)sin[(62.8m$^{-1}$)x]cos[(628s$^{-1}$)t]. Assuming $\pi$ = 3.14, what is the length of string (in m)?
\qn{2} A body of mass 1 kg is suspended by means of a spring and a thread between two walls as shown in the figure. The measure of both angles between the thread and the horizontal and the spring and the horizontal is shown in the figure. What is the initial acceleration (in m/ s$^2$) immediately after the string is cut (use g = 10 m/s$^2$)\par
\qfig{p1_phy_s3_q02_1}
\qn{3} In a photoelectric experimental setup ultraviolet light of wavelength 350 nm is incident on emitter plate (E). The work function of the emitter plate is $\phi$ = 2.2eV. AB is a uniform wire resistor having length L = 100 cm. Resistance of the wire AB is 10 $\Omega$ and the emf of the battery is V = 10 volt. The sliding contact J can be moved along the wire AB. When the sliding contact is placed at end B of the wire resistor the micro-ammeter shows a reading of i = $\pm$6$\mu$A. Assume that photoelectric current is very small compared to the current through the cell. Find the reading of the ammeter (in $\mu$A) when the slider is moved to end A of the wire.\par
\qfig{p1_phy_s3_q03_1}
\qn{4} A smooth fixed cylindrical glass lies on the ground. A small particle of mass m lies on the bottom surface touching circumference as shown in the figure. The minimum value of V$_0$ (along the circumference) such that particle just reaches at the top of glass is $\sqrt{\dfrac{\text{Kgh}}{3}}$. Find the value of K.\par
\qfig{p1_phy_s3_q04_1}
\qn{5} A block of weight 100 N is suspended by copper and steel wires of same cross sectional area 0.5 cm$^2$ and, length $\sqrt{3}$m and 1 m, respectively. Their other ends are fixed on a ceiling as shown in figure. The angles subtended by copper and steel wires with ceiling are 30$^\text{o}$ and 60$^\text{o}$ , respectively. If elongation in copper wire is ($\Delta\ell_\text{C}$) and elongation in steel wire is ($\Delta\ell_\text{S}$), then the ratio $\dfrac{\Delta\ell_\text{C}}{\Delta\ell_\text{S}}$ is\_\_\_\_.\par
[Young's modulus for copper and steel are 1x10$^{11}$N/m$^2$ and 2 $\times$10$^{11}$N/m$^2$ respectively]\par
\qfig{p1_phy_s3_q05_1}
\qn{6} A water channel of width b = 17 m lies between two frictionless platforms A and B. A plank (P) of length $\ell$ = 2 m rest near the platform A. A block of mass m = 25 kg sliding on the platform A lands on the plank and stays on it. Force $\vec{\text{F}}$ of water resistance on the plank is proportional to velocity of the plank ($\vec{\text{v}}_\text{pr}$) relative to the water and given by the law $\vec{\text{F}}$ = $-$k$\vec{\text{v}}_\text{pr}$, where k = 15N -s/m. Calculate the minimum speed u$_\text{min}$ of the block on platform A so that plank reaches the platform B. (Answer in m /s)\par
\qfig{p1_phy_s3_q06_1}
\qn{7} When a resistance R is connected in series with an element A, the electric current is found to be lagging behind the applied AC voltage V by an angle 30$^\text{o}$.When the same resistance is connected in series with element B, current leads by 60$^\text{o}$.When R, A, B are connected in series, the current now leads voltage by $\theta$ which is equal to tan$^{-1}$(K/$\sqrt{3}$), then the value of K is (assume same AC source is used in all cases)
\qn{8} The wedge shown in the figure can slide without friction on a horizontal tabletop. The mass of the wedge is M and its angle of elevation is $\alpha$ = 30$^\text{o}$. A body of mass m slides down without friction along the wedge, the path of the body makes an angle $\phi$ = 60$^\text{o}$ with the ground. Find the ratio of the masses, m/M. (Initially system is at rest)\par
\qfig{p1_phy_s3_q08_1}
\partsec{2}{CHEMISTRY}
\sect{SECTION-I (i)}
\qn{1} Analyse the following reaction sequence.\par
\qfig{p1_chem_s1_q01_1}
Which of the following compound(s) on reaction with Cl$_2$ can produce a compound X, That is used in quantitative estimation of ``Magnesium'' using carius method?
\begin{qoptions}
\opt{A}{X}
\opt{B}{W}
\opt{C}{Z}
\opt{D}{P}
\end{qoptions}
\qfullbegin
\qn{2} Consider the following compounds and decide as to which of the following statements are true?\par
\qfig{p1_chem_s1_q02_1}
\qfullsplit
\begin{qoptions}
\opt{A}{(II) gives no reaction with Na metal, however, 1 mole of (IV) on reaction with Na metal will liberate 22.4 litres of H$_2$ gas at 1atm \& 273K}
\opt{B}{(I) will give brisk effervescence on addition of NaHCO$_3$ but will not bring any change in the colour of Br$_2$ water}
\opt{C}{(III) liberates H$_2$ gas with Na metal, gives white precipitate with Tollen's reagent but does not respond towards Lucas reagent or 2, 4-DNP test.}
\opt{D}{(IV) gives turbidity with anhydrous ZnCl$_2$ + HCl}
\end{qoptions}
\qfullend
\qn{3} A 0.1 M solution of Cu$^+$ ions is stirred with an excess of potassium cyanide sufficient to convert all the Cu$^+$ ions to the corresponding cuprocyanide complex [Cu(CN)$_4$]$^{-3}$ and in addition to providing the solution with an excess of CN$^-$ equal to 0.2 M. Calculate the maximum pH of the solution when the final solution is treated with hydrogen sulphide to maintain [H$_2$S] = 0.1M and such that the precipitation of cuprous sulphide is prevented. The instability constant for [Cu(CN)$_4$]$^{-3}$ is 6.4 $\times$ 10$^{-15}$ M$^4$, K$_\text{a, overall}$ of H$_2$S = 1.6 $\times$ 10$^{-21}$ M, K$_\text{sp}$ of Cu$_2$S = 2.56 $\times$ 10$^{-27}$M$^3$.
\begin{qoptions}
\opt{A}{4.0}
\opt{B}{10.0}
\opt{C}{10.8}
\opt{D}{3.2}
\end{qoptions}
\qn{4} n containers having volumes v, 2v, 3v........nv containing gaseous moles n, 4n, 9n,........n$^3$ respectively are connected with stopcock at same temperature. If pressure of first container is P, then the final pressure when all the stopcocks are opened is
\begin{qoptions}
\opt{A}{$\dfrac{\text{n}(2\text{n}+1)\text{P}}{3}$}
\opt{B}{$\dfrac{(2\text{n}+1)\text{P}}{3\text{n}}$}
\opt{C}{$\dfrac{(2\text{n}+1)\text{P}}{6\text{n}}$}
\opt{D}{$\dfrac{(2\text{n}+1)\text{P}}{3}$}
\end{qoptions}
\sect{SECTION-I (ii)}
\qn{5} Which of the following statements is/are correct regarding solubility in water ?
\begin{qoptions}
\opt{A}{Hexan-1-ol is less soluble than cyclohexanol}
\opt{B}{4-methyl phenol is more soluble than heptan-1-ol}
\opt{C}{3-Ethyl hexanol is more soluble than octan-2-ol}
\opt{D}{Cyclooctane-1, 4-diol is more soluble than hexan-2-ol}
\end{qoptions}
\qfullbegin
\qn{6} With reference to the scheme given below, which of the following statement(s) is/are true about compounds (A), (B), (C) and (D)?\par
\qfig{p1_chem_s2_q06_1}
\qfullsplit
\begin{qoptions}
\opt{A}{P is chiral and is optically pure.}
\opt{B}{Compound Q gives a resolvable mixture on reaction with Hg(OAc)$_2$ followed by reduction with NaBH$_4$}
\opt{C}{Compound R exhibits geometrical isomerism.}
\opt{D}{S is a mixture of diastereomers.}
\end{qoptions}
\qfullend
\qn{7} Which of the following statement(s) is (are) \textbf{incorrect}?
\begin{qoptions}
\opt{A}{A green coloured solution is obtained when thiourea is tested for nitrogen in lassaigne test.}
\opt{B}{It is not necessary to carry out the chromyl chloride test in a dry test tube.}
\opt{C}{In Pb(NO$_3$)$_2$, the brown ring test can be performed with its water extract}
\opt{D}{Suspension of CdCO$_3$ gives black ppt. with sodium sulphide solution.}
\end{qoptions}
\qn{8} Which of the following statement(s) is/are \textbf{true}?
\begin{qoptions}
\opt{A}{Soluble bicarbonates give white precipitate with MgCl$_2$ in cold.}
\opt{B}{Soluble calcium bicarbonates give white precipitate with dilute ammonium solution.}
\opt{C}{All bicarbonates are generally soluble in water.}
\opt{D}{Hg(II) chloride forms a reddish brown precipitate in a solution of sodium carbonate.}
\end{qoptions}
\qn{9} A 150 mL solution of I$_2$ is divided into two unequal parts. `I' part reacts with hypo solution in acidic medium where 15 mL of 0.4 M hypo was consumed. `II' part was added to 100 mL of 0.3 M hot and concentrated NaOH solution. Residual base required 10 mL of 0.3 M H$_2$SO$_4$ solution for complete neutralization. Then
\begin{qoptions}
\opt{A}{Number of moles of I$_2$ in part-I is 3 $\times$ 10$^{-3}$}
\opt{B}{Initial concentration of I$_2$ is 0.1M}
\opt{C}{Number of moles of I$_2$ in part-II is 15 $\times$ 10$^{-3}$}
\opt{D}{Number of moles of I$_2$ in part-II is 12 $\times$ 10$^{-3}$}
\end{qoptions}
\qn{10} In a large sample of Hydrogen atoms, all atoms are in 9th excited state. Identify the correct statements when all the excited atoms reach ground state without producing any spectral line in Balmer series.
\begin{qoptions}
\opt{A}{Total 36 different spectral lines are produced}
\opt{B}{Total 37 different spectral lines are produced}
\opt{C}{8 different spectral lines are produced in Lyman series}
\opt{D}{9 different spectral lines are produced in Lyman series}
\end{qoptions}
\sect{SECTION-III}
\qn{1} The total number(s) of stable conformers with non-zero dipole moment for the following compound is/are\par
\qfig{p1_chem_s3_q01_1}
\qn{2} The total number of compounds having p$\pi$--d$\pi$ bonding among the molecules given below are\par
SO$_2$, SO$_3$, CO$_2$, H$_3$PO$_4$, H$_3$BO$_3$, HClO$_4$, HNO$_3$, HOCl
\qn{3} A solution of Hg$^{2+}$ ion on treatment with a solution of cobalt (II) thiocyanate gives rise to a deep blue crystalline precipitate. The coordination number of mercury in the deep blue coloured compound is
\qn{4} Cl$_2$ + F$_2$(excess) $\xrightarrow{\text{(573K)}}$ A $\xrightarrow{\text{Uranium}}$ B + C inter halogen compound\par
The total no. of $\ell$.p's in A is $\longrightarrow$ x\par
The total no. of $\ell$.p's in B is $\longrightarrow$ y\par
The total no. of $\ell$.p's in C is $\longrightarrow$ z\par
x + y $-$ z = 15 + b. Then `b' is \_\_\_\_\_\_\_
\qn{5} Consider the following reaction sequence.\par
\qfig{p1_chem_s3_q05_1}
The degree of unsaturation of compound E is
\qn{6} How many of the following compounds exhibit positive ``Silver mirror test'' with Tollen's reagent?\par
\qfig{p1_chem_s3_q06_1}
(i) 3, 3$-$dimethyl-1-butyne\par
(k) Acetophenone
\qn{7} From the following choose the number of correct statement(s) :-\par
(i) Stability of lyophillic colloids is mainly due to the strong interaction between dispersed phase and dispersion medium.\par
(ii) Entropy change for micelle formation in aqueous solution is positive.\par
(iii) Gelatin has considerably low value of gold number and is an effective protective agent.\par
(iv) Zeta potential is also responsible for stability of lyophobic colloids.\par
(v) Surface tension of lyophilic colloids is lesser than that of the dispersion medium.\par
(vi) A positively charged ion can cause the precipitation of As$_2$S$_3$ sol\par
(vii) Cellulose solution is an example of multimolecular colloid system.\par
(viii) The combination of two layers of same charge around the colloidal particle is called Helmholtz electrical double layer.
\qn{8} Number of crystal systems having at least 2 interfacial angles equal = x\par
Number of crystal systems having none of the interfacial angles equal = y\par
Number of crystal systems having none of the edge lengths equal = z\par
Then, x $-$ y + z is\_\_\_\_\_
\partsec{3}{MATHEMATICS}
\sect{SECTION-I (i)}
\qn{1} If the curves y = 2(x $-$ a)$^2$ and y = e$^{2x}$ touches each other, then `a' is less than-
\begin{qoptions}
\opt{A}{$-$3}
\opt{B}{$-$2}
\opt{C}{$-$1}
\opt{D}{0}
\end{qoptions}
\qn{2} If normal to the curve y = x$^2$ forms the shortest chord, then the square of slope of normal is equal to
\begin{qoptions}
\opt{A}{$\dfrac{1}{\sqrt{2}}$}
\opt{B}{$\dfrac{1}{2}$}
\opt{C}{$\sqrt{2}$}
\opt{D}{2}
\end{qoptions}
\qn{3} If I(k) = $\displaystyle\int_0^\infty\dfrac{\ell\text{nx}}{\text{x}^2 + \text{kx} + \text{k}^2}\text{dx}$ (k $>$ 0) then
\begin{qoptions}
\opt{A}{I(k) $>$ $\dfrac{\ell\text{nk}}{3\text{k}}$}
\opt{B}{I(k) $<$ $\dfrac{\ell\text{nk}}{4\text{k}}$}
\opt{C}{I(1) = 1}
\opt{D}{I(1) $<$ $\dfrac{\pi}{\sqrt{3}}$}
\end{qoptions}
\qn{4} Let z and w be two complex numbers such that $|$z$|$ $\leq$ 1, $|$w$|$ $\leq$ 1 and $|$z + i$\bar{\text{w}}|$ = $|$z $-$ i$\bar{\text{w}}|$ = 2 then which of the following is \textbf{incorrect}?
\begin{qoptions}
\opt{A}{Number of complex number z satisfying above condition are 2}
\opt{B}{Re(z) = Im(w)}
\opt{C}{Im(z) = Re(w)}
\opt{D}{$|$Z$|$ = $|$w$|$ $\neq$ 1}
\end{qoptions}
\sect{SECTION-I (ii)}
\qn{5} If $\displaystyle\lim_{n\to\infty}\left(\dfrac{1}{1^2} + \dfrac{1}{2^2} + \ldots + \dfrac{1}{n^2}\right) = \dfrac{\pi^2}{6}$ and S = $\displaystyle\lim_{n\to\infty}\left(\dfrac{1}{1^3.2^3} + \dfrac{1}{2^3.3^3} + \dfrac{1}{3^3.4^3} + \ldots + \dfrac{1}{n^3(n+1)^3}\right)$ then S + $\pi^2$ is.
\begin{qoptions}
\opt{A}{A rational number}
\opt{B}{An irrational number}
\opt{C}{An integer}
\opt{D}{Non integer}
\end{qoptions}
\qn{6} f : N $\longrightarrow$ A and g : N $\longrightarrow$B are two onto functions such that f(n) = $\left[\text{nsec}^2\theta\right]$ and g(n) = $\left[\text{n cos ec}^2\theta\right]$, for some $\theta$. If sec$^2\theta$ is irrational, then (where[.] denotes the greatest integer function \& N is set of natural numbers)
\begin{qoptions}
\opt{A}{A $\cap$ B = $\phi$}
\opt{B}{A $\cup$ B = N}
\opt{C}{f and g are bijective}
\opt{D}{A $\cap$ B is non-empty finite set}
\end{qoptions}
\qn{7} 16 players P$_1$, P$_2$, ......... P$_{16}$ play a knock out tournament. Assuming that the players are paired at random in each round. If all the players are of equal strength, then the probability that P$_1$ reaches final under the condition P$_2$ reaches semifinal but not in final is $\dfrac{\text{a}}{\text{b}}$ (where a and b are co-primes), then which of the following is true?
\begin{qoptions}
\opt{A}{$\dfrac{\text{a}}{\text{b}}$ = $\dfrac{1}{15}$}
\opt{B}{$\dfrac{\text{a}}{\text{b}}$ = $\dfrac{2}{15}$}
\opt{C}{Number of divisors of a$^2$b = 12}
\opt{D}{Sum of all devisors of a$^2$b = 24}
\end{qoptions}
\qn{8} Given $\overrightarrow{\text{A}}$(x) = x$\hat{\text{i}}$ + x$\hat{\text{j}}$ + 5$\hat{\text{k}}$, $\overrightarrow{\text{B}}$ (x) = $\hat{\text{i}}$ + x$\hat{\text{j}}$, $\overrightarrow{\text{C}}$ (x) = $\hat{\text{i}}$ + $\dfrac{1}{5}\hat{\text{j}}$ + x$\hat{\text{k}}$. Let $\alpha$, $\beta$, $\gamma$ denotes three values of x for which $\overrightarrow{\text{A}}$ (x), $\overrightarrow{\text{B}}$ (x) and $\overrightarrow{\text{C}}$ (x) are coplanar, then which of the following is correct?
\begin{qoptions}
\opt{A}{[$\alpha$] + [$\beta$] + [$\gamma$] = $-$ 1}
\opt{B}{[$\alpha$] + [$\beta$] + [$\gamma$] = 0}
\opt{C}{Global maximum value of volume of parallelpiped formed by three vectors whose conterminous edges are $\overrightarrow{\text{A}}$ (x), $\overrightarrow{\text{B}}$ (x) \& $\overrightarrow{\text{C}}$ (x) is 4 for x $\in$ [$-$2, 3]}
\opt{D}{(1 + $\alpha$)(1 + $\beta$)(1 + $\gamma$) = $-$4 (where [.] denotes greatest integer function)}
\end{qoptions}
\qn{9} A triangle with vertices P(2, 5), Q(5, 2) and R($-$1, $-$1) is inscribed in the rectangular hyperbola xy $-$ x $-$ y $-$ 3 = 0 and let H $\equiv$ ($\lambda$, $\lambda$) is orthocenter of $\Delta$PQR, let A, B and C be the feet of perpendicular from P, Q and R on QR, PR and PQ respectively and r is in radius of $\Delta$ABC, then
\begin{qoptions}
\opt{A}{Circumcentre of $\Delta$PQR is $\left(-\dfrac{3}{2}, -\dfrac{3}{2}\right)$}
\opt{B}{Incentre of $\Delta$ABC is (3, 3)}
\opt{C}{Equation of AB is x + y $-$ $\dfrac{26}{5}$ = 0}
\opt{D}{r = $\dfrac{2\sqrt{2}}{5}$}
\end{qoptions}
\qn{10} Let f(x) = $\ell$n(2x $-$ x$^2$) + sin $\dfrac{\pi\text{x}}{2}$, then
\begin{qoptions}
\opt{A}{Graph of f is symmetric about the line x = 1}
\opt{B}{Graph of f is symmetric about the line x = 2}
\opt{C}{Local maximum value of f is 1}
\opt{D}{Local minimum value of f does not exist}
\end{qoptions}
\sect{SECTION-III}
\qn{1} Find sum of all integral values of r such that all roots of the polynomial\par
f(x) = x$^3$ $-$ (r $-$ 1)x$^2$ $-$11x + 4r are also integer
\qn{2} Let I = $\displaystyle\int_0^{\pi/2}\dfrac{\cos x + 4}{3\sin x + 4\cos x + 25}\text{dx}$ and J = $\displaystyle\int_0^{\pi/2}\dfrac{\sin x + 3}{3\sin x + 4\cos x + 25}\text{dx}$. If (11I + 2J) = a$\pi$ + bln $\dfrac{\text{c}}{\text{d}}$ where a, b, c and d $\in$ N and $\dfrac{\text{c}}{\text{d}}$ is not a perfect square of a rational then find the value of (a + b + c + d $-$ 55).\par
(Where N is is set of natural numbers)
\qn{3} Let a vector $\bar{\text{V}}_1$ = $\overline{\text{OP}}$, where O is origin and P is (2, 0,2). Vector $\bar{\text{V}}_1$ is rotated about origin by an angle of $\dfrac{\pi}{3}$ in a plane which is perpendicular to XOZ plane to get vector $\bar{\text{V}}_2$ where $\bar{\text{V}}_2$. $\hat{\text{j}}$ $>$ 0. Again $\bar{\text{V}}_1$ is rotated about origin by right angle in XOZ plane to get vector $\bar{\text{V}}_3$. If V is volume of tetrahedron whose coterminous edges are $\bar{\text{V}}_1$, $\bar{\text{V}}_2$ and $\bar{\text{V}}_3$ then value of $\dfrac{3\text{V}}{\sqrt{6}}$ is.
\qn{4} Suppose that f (x) is a differentiable invertible function with f(x) $\neq$ 0 and h (x) = $\displaystyle\int_1^{x}$ f(t)dt.Given that f(1) = f$'$(1) = 1 and g(x) is inverse of f(x). Let G(x) = x$^2$g(x) $-$ xh (g(x)) $\forall$x $\in$ R then G$''$(1)=
\qn{5} Let y = f(x) (f : R $\longrightarrow$ R) be an explicit function defined by the implicit equation x$^3$ + y$^3$ + 3(x$^2$ +y$^2$) + 3 (x + y) = 14 and g be the inverse function of f.\par
If $\dfrac{\text{d}}{\text{dx}}$ [f(x +g(x)). g(x +f(x))] at x = $-$ 1 + $\sqrt[3]{15}$ is equal $\lambda$(15)$^{\frac{2}{3}}$, where $\lambda$ $\in$ I, then find the value of $|\lambda|$.
\qn{6} If $<$a, b, c$>$ are direction cosines of line of intersection of three planes x + y $-$ z = 1, ( m $-$ 1)x + (3m $-$ 7) y $-$ 5z = $-$3 and ( m $-$ 2)x + (2m $-$ 5) y = $-$4 for some m $\in$ R, then value of $\dfrac{\text{a}+\text{b}}{\text{c}}$ is \_\_\_\_\_.
\qn{7} Consider a complex number z = $\dfrac{\sqrt{5}+1}{4}$ + i$\left(\dfrac{\sqrt{5}-1}{2}\right)\cos\dfrac{\pi}{10}$ and a set S = \{Z$_\text{p}$$|$Z$_\text{p}$ = Z$^\text{p}$, p = 0, 1, 2, 3,. . . .9\} Number of ordered pairs (z$_\text{i}$, z$_\text{j}$) where z$_\text{i}$, z$_\text{j}$ $\in$ S such that arg (z$_\text{i}$) $-$ arg (z$_\text{j}$) = $\dfrac{2\pi}{5}$ (where arg z $\in$ [0,2$\pi$) is.
\qn{8} A curve y = f(x) is passing through (2, 0 ) and slope of its tangent at any point (x, y) is $\dfrac{(\text{x}+1)^2 + \text{y} - 3}{\text{x}+1}$. Then the area bounded by the curve y = f(x) and x $-$ axis is $\dfrac{\text{m}}{\text{n}}$ where H.C. F of m \& n is 1. Then value of m + n is.
\end{multicols}
\end{document}
